% Generated draft: requires mathematical and source review.
\documentclass[11pt]{article}
\usepackage[margin=25mm]{geometry}
\usepackage{fontspec}
\setmainfont{Latin Modern Roman}
\usepackage{amsmath,amssymb,mathtools,unicode-math}
\setmathfont{Latin Modern Math}
\usepackage{enumitem,needspace,xurl,hyperref}
\setlength{\parindent}{0pt}
\setlength{\parskip}{5pt}
\setlength{\emergencystretch}{3em}
\newcommand{\N}{\mathbb{N}}
\newcommand{\Z}{\mathbb{Z}}
\newcommand{\Q}{\mathbb{Q}}
\newcommand{\R}{\mathbb{R}}
\newcommand{\C}{\mathbb{C}}
\newcommand{\F}{\mathbb{F}}
\newcommand{\A}{\mathbb{A}}
\newcommand{\PP}{\mathbb{P}}
\newcommand{\E}{\mathbb{E}}
\newcommand{\eps}{\varepsilon}
\newcommand{\dd}{\,\mathrm{d}}
\begin{document}
\section*{UnsolvedMath ID: 30001998}
\textbf{Source problem code:} catalogue entry
\subsection{Definitions and mathematical statement}
Let \(R\) be a discrete valuation ring, let \(S=\operatorname{Spec}R\), and let \(f\colon G\to S\) be a group scheme. Write \(\mathcal O_G\) for its structure sheaf and define the graded cohomology \(R\)-algebra
\[
H^\bullet(G,\mathcal O_G)=\bigoplus_{i\geq 0}H^i(G,\mathcal O_G),
\qquad
H^i(G,\mathcal O_G)\otimes_R H^j(G,\mathcal O_G)\xrightarrow{\smile}H^{i+j}(G,\mathcal O_G),
\]
where multiplication is the cup product induced by multiplication on \(\mathcal O_G\), and the scalar action comes from \(f\).

The supplied problem asks for a description of \(H^\bullet(G,\mathcal O_G)\) for group schemes over discrete valuation rings, and for a classification of anti-affine group schemes over such rings.

A possible relative meaning of anti-affineness is the condition that the structural map
\[
R\longrightarrow\Gamma(G,\mathcal O_G)
\]
is an isomorphism. Under this interpretation, the second request concerns group schemes satisfying this condition. The supplied record does not define anti-affineness over \(R\), however, so this interpretation requires confirmation; no condition on geometric fibers is silently imposed.

The intended hypotheses on \(G\) are unspecified: the record gives no finite type, smoothness, flatness, separatedness, connectedness, or commutativity assumptions. It also gives no restrictions on the discrete valuation ring beyond that designation, and does not specify the equivalence relation or invariants intended for the classification. Thus these requests describe a research problem with unresolved scope, rather than a fully specified classification theorem. The displayed sheaf cohomology algebra is defined without adding these hypotheses; calling it coherent cohomology does not supply missing finiteness assumptions.
\subsection{Short English statement}
Describe the graded structure-sheaf cohomology ring of group schemes over a discrete valuation ring, and classify the anti-affine group schemes in that setting. The allowed group schemes, the precise relative meaning of anti-affineness, and the classification criterion require clarification.
\subsection{Sources}
\begin{itemize}
\item Supplied local UnsolvedMath registry record, permanent ID \(30001998\), title \emph{Cohomology and Classification of Group Schemes over DVRs}, source problem code \texttt{catalogue entry}.
\item Supplied report citation: \emph{Enveloping Algebras and Geometric Representation Theory} (2012), Oberwolfach Reports, pp.\ 733--809, DOI \texttt{10.4171/owr/2012/13}. Source and supplied literature-source link, labeled ``Oberwolfach Report 13 2012'': \url{https://doi.org/10.4171/owr/2012/13}. The record attributes the problem to workshop open-problem discussions documented in Oberwolfach Reports; it names no individual proposer and gives the proposed year as 2012.
\item Literature assessment dated exactly \texttt{2026-08-21}: the supplied curation labels the problem open and states that the coherent cohomology ring over a discrete valuation ring and the structure of anti-affine group schemes remain unresolved in the generality stated. This is supplied dated literature triage derived from the OWR curation, not an independently verified assessment. Time-sensitive status and cited claims require independent verification.
\item The supplied statement status is \texttt{corrected\_verified}, with a recorded manual semantic, typographic, and LaTeX review describing the cleanup as meaning-preserving and self-contained. This supplied editorial verification does not establish mathematical correctness or independently verify the literature status.
\end{itemize}
\end{document}
